% ============================================================================
% Chapter 28: Introduction to MS Excel
% ============================================================================
\chapter{Introduction to MS Excel}
\label{ch:excel-intro}

\begin{objectives}
  \item Open MS Excel and identify parts of the Excel window
  \item Understand cells, rows, columns, and cell addresses
  \item Enter data into cells (text, numbers, dates)
  \item Navigate a worksheet using keyboard and mouse
\end{objectives}

\bytetip[tip]{Excel is like a super-powered calculator combined with a smart notebook. It can do maths, organise lists, create charts, and much more. Let's explore!}

\section{The Excel Window}
\label{sec:excel-window}
\index{MS Excel!interface}

% --- Figure 28.1: Excel Window Diagram ---
\begin{figure}[H]
\centering
\begin{tikzpicture}[scale=0.7]
  \draw[NavyBlue, line width=2pt, rounded corners=3pt] (0,0) rectangle (16,11);
  
  % Title bar
  \fill[NavyBlue!80] (0.1,10.2) rectangle (15.9,11);
  \node[text=white, font=\tiny\sffamily] at (8,10.6) {Book1 --- Microsoft Excel};
  
  % Ribbon area
  \fill[LabGreen!10] (0.1,8.8) rectangle (15.9,10.2);
  \node[font=\tiny\sffamily, text=NavyBlue] at (8,9.5) {Home \quad Insert \quad Page Layout \quad Formulas \quad Data \quad Review \quad View};
  
  % Name Box + Formula Bar
  \fill[ShadeGray] (0.1,8.0) rectangle (15.9,8.8);
  \node[draw=NavyBlue!30, fill=white, font=\tiny\sffamily\bfseries, minimum width=1.2cm, minimum height=0.5cm] at (1.2,8.4) {A1};
  \node[font=\tiny\sffamily, text=NavyBlue!50, anchor=west] at (3,8.4) {\textit{fx}\quad =SUM(A1:A5)};
  
  % Column headers
  \fill[ModuleBlue!20] (1.5,7.2) rectangle (15.5,8.0);
  \foreach \x/\c in {2.5/A, 4.5/B, 6.5/C, 8.5/D, 10.5/E, 12.5/F, 14.5/G}{
    \node[font=\tiny\sffamily\bfseries, text=ModuleBlue] at (\x,7.6) {\c};
  }
  
  % Row numbers
  \fill[NavyBlue!10] (0.1,1.0) rectangle (1.5,7.2);
  \foreach \y/\r in {6.5/1, 5.5/2, 4.5/3, 3.5/4, 2.5/5, 1.5/6}{
    \node[font=\tiny\sffamily\bfseries, text=NavyBlue] at (0.8,\y) {\r};
  }
  
  % Cell grid
  \fill[white] (1.5,1.0) rectangle (15.5,7.2);
  \foreach \x in {1.5,3.5,...,15.5}{
    \draw[NavyBlue!15, line width=0.3pt] (\x,1.0) -- (\x,7.2);
  }
  \foreach \y in {1.0,2.0,...,7.0}{
    \draw[NavyBlue!15, line width=0.3pt] (1.5,\y) -- (15.5,\y);
  }
  
  % Highlighted cell
  \draw[ModuleBlue, line width=1.5pt] (1.5,6.2) rectangle (3.5,7.2);
  
  % Sheet tabs
  \fill[ShadeGray] (0.1,0) rectangle (15.9,1.0);
  \node[draw=NavyBlue!30, fill=white, font=\tiny\sffamily, rounded corners=2pt, minimum width=1.5cm, minimum height=0.4cm] at (2,0.5) {Sheet1};
  \node[draw=NavyBlue!15, fill=ShadeGray, font=\tiny\sffamily, rounded corners=2pt, minimum width=1.5cm, minimum height=0.4cm] at (4,0.5) {Sheet2};
  
  % Labels
  \draw[FunCoral, line width=1pt, ->] (-1.5,8.4) -- (0,8.4);
  \node[left, font=\scriptsize\sffamily\bfseries, text=FunCoral, align=right] at (-1.5,8.4) {Name Box \&\\Formula Bar};
  
  \draw[ModuleBlue, line width=1pt, ->] (-1.5,7.6) -- (1.4,7.6);
  \node[left, font=\scriptsize\sffamily\bfseries, text=ModuleBlue, align=center] at (-1.5,7.6) {Column\\Headers};
  
  \draw[NavyBlue, line width=1pt, ->] (-1.5,4.5) -- (0,4.5);
  \node[left, font=\scriptsize\sffamily\bfseries, text=NavyBlue, align=center] at (-1.5,4.5) {Row\\Numbers};
  
  \draw[FunSky, line width=1pt, ->] (-1.5,0.5) -- (0,0.5);
  \node[left, font=\scriptsize\sffamily\bfseries, text=FunSky] at (-1.5,0.5) {Sheet Tabs};
\end{tikzpicture}
\caption{The MS Excel window with all key elements labelled}
\label{fig:excel-window}
\end{figure}

\section{Understanding Cell Addresses}
\label{sec:cell-address}
\index{cell address}\index{MS Excel!cells}

% --- Figure 28.2: Cell Address Explained ---
\begin{figure}[H]
\centering
\begin{tikzpicture}
  % 4x4 grid
  \foreach \x in {0,2,4,6}{
    \draw[NavyBlue!30, line width=0.5pt] (\x,0) -- (\x,4);
  }
  \foreach \y in {0,1,2,3,4}{
    \draw[NavyBlue!30, line width=0.5pt] (0,\y) -- (6,\y);
  }
  
  % Headers
  \foreach \x/\c in {1/A, 3/B, 5/C}{
    \node[font=\small\sffamily\bfseries, text=ModuleBlue] at (\x,4.4) {\c};
  }
  \foreach \y/\r in {3.5/1, 2.5/2, 1.5/3, 0.5/4}{
    \node[font=\small\sffamily\bfseries, text=NavyBlue] at (-0.4,\y) {\r};
  }
  
  % Highlight C3
  \fill[FunYellow!40] (4,1) rectangle (6,2);
  \draw[FunYellow!80!black, line width=2pt] (4,1) rectangle (6,2);
  \node[font=\small\sffamily\bfseries] at (5,1.5) {42};
  
  % Column arrow
  \draw[ModuleBlue, line width=1.5pt, ->] (5,4.4) -- (5,4.8);
  \node[above, font=\small\sffamily\bfseries, text=ModuleBlue] at (5,4.8) {Column C};
  
  % Row arrow
  \draw[NavyBlue, line width=1.5pt, ->] (-0.4,1.5) -- (-0.8,1.5);
  \node[left, font=\small\sffamily\bfseries, text=NavyBlue] at (-0.8,1.5) {Row 3};
  
  % Speech bubble
  \node[draw=FunYellow!80!black, fill=FunYellow!10, rounded corners=6pt,
        font=\small\sffamily\bfseries, text=NavyBlue, inner sep=5pt] 
    at (9,1.5) {My address is \textbf{C3}!};
  \draw[FunYellow!80!black, line width=1pt, ->] (7.5,1.5) -- (6.1,1.5);
\end{tikzpicture}
\caption{Cell C3 --- column letter + row number = cell address}
\label{fig:cell-address}
\end{figure}

\begin{theorybox}[Cell Address]
Every cell in Excel has a unique address made of its \textbf{column letter} and \textbf{row number}. For example, the cell in column C, row 3 has the address \textbf{C3}. This is how Excel knows which cell you're talking about in formulas.
\end{theorybox}

\bytetip[fun]{An Excel worksheet has \textbf{16,384 columns} (A to XFD) and \textbf{1,048,576 rows}. That's over \textbf{17 billion cells} in a single sheet! You'd need to live several lifetimes to fill them all.}

\begin{funfact}
The first spreadsheet program was called \textbf{VisiCalc}, released in 1979 for the Apple II. It was so popular that people bought Apple II computers \emph{just} to run VisiCalc. It's considered one of the most important software programs ever created!
\end{funfact}

\begin{reviewqs}
  \item Name five parts of the Excel window.
  \item What is a cell address? Give three examples.
  \item How many columns and rows does an Excel worksheet have?
  \item Where is the Formula Bar, and what does it show?
  \item What are sheet tabs, and what do they do?
\end{reviewqs}

\begin{challenge}
Open Excel and type data for 5 classmates: Name (column A), Age (column B), Favourite Subject (column C). Bold the header row. Adjust column widths so all text is visible. Save as ``ClassData.xlsx.''
\end{challenge}
